\subsection{INTEGRATION OF COMPARISON RELATIONS: I. WEAK RELATIONS}

PROPOSITION 1. Let \(\mathbf{f}\) be a regulated vector function, \(g\) a regulated function \(\geqslant 0\) on an interval \([a, +\infty[\) , such that \(\int_{a}^{+\infty} g(t) dt > 0\) . The relation \(\mathbf{f} \preccurlyeq g\) as \(x\) tends to \(+\infty\) implies that \(\int_{a}^{x} \mathbf{f}(t) dt \preccurlyeq \int_{a}^{x} g(t) dt\) . If the integral \(\int_{a}^{+\infty} g(t) dt\) converges then the integral \(\int_{a}^{+\infty} \mathbf{f}(t) dt\) converges absolutely.

Indeed, by hypothesis there exist a \(b \geqslant a\) and a number \(c' > 0\) such that

\[
\| \mathbf {f} (x) \| \leqslant c ^ {\prime} g (x) \quad \text {   for   } x \geqslant b,
\]

whence

\[
\left\| \int_ {b} ^ {x} \mathbf {f} (t) d t \right\| \leqslant \int_ {b} ^ {x} \| \mathbf {f} (t) \| d t \leqslant c ^ {\prime} \int_ {b} ^ {x} g (t) d t;
\]

since, moreover, one may suppose \(b\) so large that \(\int_{a}^{b} g(t) dt > 0\) , there exists a \(c'' > 0\) such that \(\left\| \int_{a}^{b} \mathbf{f}(t) dt \right\| \leqslant c'' \int_{a}^{b} g(t) dt\) ; on putting \(c = \max(c', c'')\) one thus has

\[
\left\| \int_ {a} ^ {x} \mathbf {f} (t) d t \right\| \leqslant c \int_ {a} ^ {x} g (t) d t
\]

for every \(x \geqslant b\) , whence the proposition.

COROLLARY 1. If \(f\) and \(g\) are regulated functions and \(\geqslant 0\) on the interval \([a, +\infty[\) , and such that \(f \succcurlyeq g\) , and if \(\int_{a}^{+\infty} g(t) dt = +\infty\) , then \(\int_{a}^{+\infty} f(t) dt = +\infty\) .

COROLLARY 2. If \(f\) and \(g\) are \(\geqslant 0\) and not identically zero on \([a, +\infty[\) , and such that \(f \asymp g\) , then \(\int_{a}^{x} f(t) dt \asymp \int_{a}^{x} g(t) dt\) .

